% GENERATED FILE. Edit canonical lessons, then run npm run generate:books.
% canonical-source-hash: fnv1a64:9a89af15d7ab98a8
% canonical-lessons: SA-C348-sosanam, SA-C348-nyayyah, SA-C348-vidhih, SA-C348-paksapatah, SA-C348-ullanghayati

\chapter{Fair and Unfair}
\label{ch:sa-fair-and-unfair}
\hlchaptermodality{348}

\begin{chapteropening}
\textbf{By the end of this chapter:} \emph{I can say whether people are treated fairly.}

Complete the last lesson of chapter 348: I can say whether people are treated fairly.
\end{chapteropening}

\section[\texorpdfstring{\sk{शोषणम्} (śoṣaṇam)}{शोषणम् (śoṣaṇam)}]{\sk{शोषणम्} (śoṣaṇam) — exploitation}

\label{lesson:SA-C348-sosanam}



\begin{quote}
Before the new one: say the Sanskrit for makes (someone) understand, then the Sanskrit for what someone means.
\end{quote}

\subsection*{You'll want to know: \sk{शोषणम्}}
\textbf{\sk{शोषणम्}} — \emph{śoṣaṇam} — \textquotedblleft{}exploitation\textquotedblright{}.

\begin{quote}
\sk{धनिकः} \sk{कर्मकरस्य} \sk{शोषणं} \sk{करोति।}
\end{quote}

\begin{quote}
\emph{dhanikaḥ karmakarasya śoṣaṇaṃ karoti.}
\end{quote}

\begin{quote}
The rich man exploits the worker.
\end{quote}

A neuter noun.

\begin{cousinweb}
\textbf{\sk{शोषणम्}} (\emph{śoṣaṇam}) comes from the root √śuṣ, \textquotedblleft{}to dry\textquotedblright{}, in its causative form, \textquotedblleft{}to make dry\textquotedblright{}; MW gives \emph{śoṣaṇa} \textquotedblleft{}drying up, draining, suction\textquotedblright{}, and \textquotedblleft{}exploitation\textquotedblright{} is the sense modern Sanskrit gives it (MW).
\end{cousinweb}

\begin{grammarlens}[title={What you've built}]
One more word for giving reasons and opinions.
\end{grammarlens}

\subsection*{Your turn}
\begin{itemize}
\raggedright
  \item \emph{Say it:} \emph{śoṣaṇam}
  \item \emph{Say it:} \emph{śoṣaṇam}, once more
  \item \emph{Say it:} \emph{dhanikaḥ karmakarasya śoṣaṇaṃ karoti.} — The rich man exploits the worker.
  \item \emph{Recall:} say the Sanskrit for makes (someone) understand, then the Sanskrit for what someone means, then say \emph{śoṣaṇam} again
\end{itemize}

\begin{tcolorbox}[breakable,colback=teal!4,colframe=teal!35!black,title={Before you move on}]
What does \sk{शोषणम्} mean? (\textquotedblleft{}Exploitation\textquotedblright{}.) Say it once more.
\end{tcolorbox}
\section[\texorpdfstring{\sk{न्याय्यः} (nyāyyaḥ)}{न्याय्यः (nyāyyaḥ)}]{\sk{न्याय्यः} (nyāyyaḥ) — fair; right}

\label{lesson:SA-C348-nyayyah}



\begin{quote}
Before the new one: say the Sanskrit for what someone means, then the Sanskrit for exploitation.
\end{quote}

\subsection*{You'll want to know: \sk{न्याय्यः}}
\textbf{\sk{न्याय्यः}} — \emph{nyāyyaḥ} — \textquotedblleft{}fair; right\textquotedblright{}.

\begin{quote}
\sk{सः} \sk{निर्णयः} \sk{न्याय्यः} \sk{न} \sk{अस्ति।}
\end{quote}

\begin{quote}
\emph{saḥ nirṇayaḥ nyāyyaḥ na asti.}
\end{quote}

\begin{quote}
That decision is not fair.
\end{quote}

\begin{cousinweb}
\textbf{\sk{न्याय्यः}} (\emph{nyāyyaḥ}) is built on \emph{nyāya}, \textquotedblleft{}rule, right manner\textquotedblright{}; MW gives \emph{nyāyya} \textquotedblleft{}regular, correct, right, fit, proper\textquotedblright{} (MW).
\end{cousinweb}

\begin{grammarlens}[title={What you've built}]
One more word for giving reasons and opinions.
\end{grammarlens}

\subsection*{Your turn}
\begin{itemize}
\raggedright
  \item \emph{Say it:} \emph{nyāyyaḥ}
  \item \emph{Say it:} \emph{nyāyyaḥ}, once more
  \item \emph{Say it:} \emph{saḥ nirṇayaḥ nyāyyaḥ na asti.} — That decision is not fair.
  \item \emph{Recall:} say the Sanskrit for what someone means, then the Sanskrit for exploitation, then say \emph{nyāyyaḥ} again
\end{itemize}

\begin{tcolorbox}[breakable,colback=teal!4,colframe=teal!35!black,title={Before you move on}]
What does \sk{न्याय्यः} mean? (\textquotedblleft{}Fair; right\textquotedblright{}.) Say it once more.
\end{tcolorbox}
\section[\texorpdfstring{\sk{विधिः} (vidhiḥ)}{विधिः (vidhiḥ)}]{\sk{विधिः} (vidhiḥ) — a precept; a rule laid down}

\label{lesson:SA-C348-vidhih}



\begin{quote}
Before the new one: say the Sanskrit for exploitation, then the Sanskrit for fair.
\end{quote}

\subsection*{You'll want to know: \sk{विधिः}}
\textbf{\sk{विधिः}} — \emph{vidhiḥ} — \textquotedblleft{}a precept; a rule laid down\textquotedblright{}.

\begin{quote}
\sk{जनः} \sk{विधिं} \sk{पालयेत्।}
\end{quote}

\begin{quote}
\emph{janaḥ vidhiṃ pālayet.}
\end{quote}

\begin{quote}
One should keep to the precept.
\end{quote}

A masculine noun.

\begin{cousinweb}
\textbf{\sk{विधिः}} (\emph{vidhiḥ}) is built on \emph{vi} and the root √dhā, \textquotedblleft{}to put, to place\textquotedblright{}; MW gives \emph{vidhi} \textquotedblleft{}a rule, formula, injunction, precept, law\textquotedblright{} (MW).
\end{cousinweb}

\begin{grammarlens}[title={What you've built}]
One more word for giving reasons and opinions.
\end{grammarlens}

\subsection*{Your turn}
\begin{itemize}
\raggedright
  \item \emph{Say it:} \emph{vidhiḥ}
  \item \emph{Say it:} \emph{vidhiḥ}, once more
  \item \emph{Say it:} \emph{janaḥ vidhiṃ pālayet.} — One should keep to the precept.
  \item \emph{Recall:} say the Sanskrit for exploitation, then the Sanskrit for fair, then say \emph{vidhiḥ} again
\end{itemize}

\begin{tcolorbox}[breakable,colback=teal!4,colframe=teal!35!black,title={Before you move on}]
What does \sk{विधिः} mean? (\textquotedblleft{}A precept; a rule laid down\textquotedblright{}.) Say it once more.
\end{tcolorbox}
\section[\texorpdfstring{\sk{पक्षपातः} (pakṣapātaḥ)}{पक्षपातः (pakṣapātaḥ)}]{\sk{पक्षपातः} (pakṣapātaḥ) — partiality; bias}

\label{lesson:SA-C348-paksapatah}



\begin{quote}
Before the new one: say the Sanskrit for fair, then the Sanskrit for a precept.
\end{quote}

\subsection*{You'll want to know: \sk{पक्षपातः}}
\textbf{\sk{पक्षपातः}} — \emph{pakṣapātaḥ} — \textquotedblleft{}partiality; bias\textquotedblright{}.

\begin{quote}
\sk{न्यायाधीशः} \sk{पक्षपातं} \sk{त्यजेत्।}
\end{quote}

\begin{quote}
\emph{nyāyādhīśaḥ pakṣapātaṃ tyajet.}
\end{quote}

\begin{quote}
A judge should put aside partiality.
\end{quote}

A masculine noun.

\begin{cousinweb}
\textbf{\sk{पक्षपातः}} (\emph{pakṣapātaḥ}) is built from \emph{pakṣa}, \textquotedblleft{}a side\textquotedblright{}, and \emph{pāta}, \textquotedblleft{}a falling\textquotedblright{}; MW gives \emph{pakṣa-pāta} \textquotedblleft{}adopting a side, siding with, partiality\textquotedblright{} (MW).
\end{cousinweb}

\begin{grammarlens}[title={What you've built}]
One more word for giving reasons and opinions.
\end{grammarlens}

\subsection*{Your turn}
\begin{itemize}
\raggedright
  \item \emph{Say it:} \emph{pakṣapātaḥ}
  \item \emph{Say it:} \emph{pakṣapātaḥ}, once more
  \item \emph{Say it:} \emph{nyāyādhīśaḥ pakṣapātaṃ tyajet.} — A judge should put aside partiality.
  \item \emph{Recall:} say the Sanskrit for fair, then the Sanskrit for a precept, then say \emph{pakṣapātaḥ} again
\end{itemize}

\begin{tcolorbox}[breakable,colback=teal!4,colframe=teal!35!black,title={Before you move on}]
What does \sk{पक्षपातः} mean? (\textquotedblleft{}Partiality; bias\textquotedblright{}.) Say it once more.
\end{tcolorbox}
\section[\texorpdfstring{\sk{उल्लंघयति} (ullaṅghayati)}{उल्लंघयति (ullaṅghayati)}]{\sk{उल्लंघयति} (ullaṅghayati) — breaks (a rule)}

\label{lesson:SA-C348-ullanghayati}



\begin{quote}
Before the new one: say the Sanskrit for a precept, then the Sanskrit for partiality.
\end{quote}

\subsection*{You'll want to know: \sk{उल्लंघयति}}
\textbf{\sk{उल्लंघयति}} — \emph{ullaṅghayati} — \textquotedblleft{}breaks (a rule)\textquotedblright{}.

\begin{quote}
\sk{यः} \sk{नियमम्} \sk{उल्लंघयति}, \sk{सः} \sk{दण्डं} \sk{प्राप्नोति।}
\end{quote}

\begin{quote}
\emph{yaḥ niyamam ullaṅghayati, saḥ daṇḍaṃ prāpnoti.}
\end{quote}

\begin{quote}
Whoever breaks a rule is punished.
\end{quote}

Past (he …): \sk{सः} \sk{नियमम्} \sk{उल्लंघितवान्}, \emph{saḥ niyamam ullaṅghitavān}.

\begin{cousinweb}
\textbf{\sk{उल्लंघयति}} (\emph{ullaṅghayati}) is built on \emph{ut}, \textquotedblleft{}up, over\textquotedblright{}, and the root √laṅgh, \textquotedblleft{}to leap\textquotedblright{}, in its causative form; MW gives the senses \textquotedblleft{}to leap over, overstep\textquotedblright{} and \textquotedblleft{}to transgress, violate\textquotedblright{} (MW).
\end{cousinweb}

\begin{grammarlens}[title={What you've built}]
One more word for giving reasons and opinions.
\end{grammarlens}

\subsection*{Your turn}
\begin{itemize}
\raggedright
  \item \emph{Say it:} \emph{ullaṅghayati}
  \item \emph{Say it:} \emph{ullaṅghayati}, once more
  \item \emph{Say it:} \emph{yaḥ niyamam ullaṅghayati, saḥ daṇḍaṃ prāpnoti.} — Whoever breaks a rule is punished.
  \item \emph{Recall:} say the Sanskrit for a precept, then the Sanskrit for partiality, then say \emph{ullaṅghayati} again
\end{itemize}

\begin{tcolorbox}[breakable,colback=teal!4,colframe=teal!35!black,title={Before you move on}]
What does \sk{उल्लंघयति} mean? (\textquotedblleft{}Breaks (a rule)\textquotedblright{}.) Say it once more.
\end{tcolorbox}
